\chapter{Theory}
In this chapter, a physical model of the underlying system will be presented
and shown to be equivalent to a discrete-time AR(1)-HMM.
A overview of the theory behind the models and inference
algorithms will be presented as well. 


\section{The Diffusion Process \& and it's properties}
This subsection present the underlying physical model of the $\text{CO}_2$ dynamics of a room 
and show that it can be described as a discrete-time AR(1) process and how the production rate of the occupants in the room 
can be modelled as a function of the occupancy state of the room. 
\subsection{Diffusion Process (Ornstein-Uhlenbeck)}
The underlying physical model of the $\text{CO}_2$ dynamics of a room can 
be described by a diffusion process. The diffusion process is a 
continuous-time stochastic process that describes the change of $\text{CO}_2$ in a room. 
\\
Let $\tilde{C(t)} = C(t) - \min_{\forall t}(C(t))$ such that the process
reverts back to a baseline of 0. The diffusion process can then be desbribed as
\begin{equation}
    \Delta \tilde{C(t)} = (P(t) - \alpha \tilde{C(t)}) \Delta t + \sigma \Delta W(t),
\end{equation}
where $P(t)$ is the $\text{CO}_2$ production rate of the occupants in the room, 
$\alpha$ is the decay rate of the $\text{CO}_2$ in the room, $\sigma$ is the diffusion coefficient and $W(t)$ is a Wiener process, which is 
interpretated as the noise in the sensors.\\
This models tells us, that if the production rate of the occupants at time $t$ is higher than 
the decay rate (or air change rate) of the room, then the $\text{CO}_2$ levels in the room will increase. 
The argument for this physical model that $\text{CO}_2$ accumate and that it
does not decay by itself [\textcolor{red}{MISSING REFERENCE}], unless it's 
ventilated out of the room. \\

\bigskip
\subsection{Proof of AR(1) process} 
We now want to show, that the underlying physical model of the $\text{CO}_2$ dynamics of a room can be described as a discrete-time AR(1) process. \\
Firstly, we recognize the diffusion process as a Ornstein-Uhlenbeck, and the solution
of this process is given by \textcite[Sec.~1.13.1]{iacus2008} as 
\begin{subequations}\label{eq:ou-solution}
\begin{align}
\tilde{C}(t) &= \frac{P}{\alpha}
      + \left( x_0 - \frac{P}{\alpha} \right) e^{-\alpha t}
      + \theta_3 \int_0^t e^{-\alpha (t-u)} \, \mathrm{d}W_u
      \label{eq:ou-param1} \\
\intertext{or, in the second parametrization,}
\tilde{C}(t) &= \mu + (x_0 - \mu) e^{-\alpha t}
      + \sigma \int_0^t e^{-\alpha (t-u)} \, \mathrm{d}W_u ,
      \qquad \mu = \frac{P}{\alpha},
      \label{eq:ou-param2}
\end{align}
\end{subequations}
where $x_0$ is the initial value of the process at time $t = 0$.

\bigskip

If we then discretize the process with a time step $\Delta t$, we can write the process as
\begin{subequations}\label{eq:ar1-proof}
\begin{align}
\tilde{C}_{t+\Delta t} &= \mu + (\tilde{C}_0 - \mu) e^{-\alpha (t + \Delta t)}
      + \sigma \int_0^{t + \Delta t} e^{-\alpha (t + \Delta t - u)} \, \mathrm{d}W_u ,
      \label{eq:ar1-full} \\
\intertext{where we can split the integral into two parts as}
\sigma \int_0^{t + \Delta t} e^{-\alpha (t + \Delta t - u)} \, \mathrm{d}W_u
  &= \sigma \biggl( e^{-\alpha \Delta t} \int_0^{t} e^{-\alpha (t-u)} \, \mathrm{d}W_u \notag \\
  &\qquad + \int_t^{t+\Delta t} e^{-\alpha (t+\Delta t-u)} \, \mathrm{d}W_u \biggr) .
      \label{eq:ar1-split} \\
\intertext{We can then write out the process as}
\tilde{C}_{t+\Delta t} &= \mu(1 - e^{-\alpha \Delta t}) + \mu e^{-\alpha \Delta t}
      + (\tilde{C}_0 - \mu) e^{-\alpha t} e^{-\alpha \Delta t} \notag \\
  \quad + \sigma e^{-\alpha \Delta t} \int_0^{t} &e^{-\alpha (t-u)} \, \mathrm{d}W_u
      + \sigma \int_t^{t+\Delta t} e^{-\alpha (t+\Delta t-u)} \, \mathrm{d}W_u
      \label{eq:ar1-expanded}
\intertext{We see then that}
\tilde{C}_{t+\Delta t} &= \biggl( \mu + (\tilde{C}_0 - \mu) e^{-\alpha t}
      + \sigma \int_0^{t} e^{-\alpha (t-u)} \, \mathrm{d}W_u \biggr) e^{-\alpha \Delta t} \notag \\
  &\quad + \mu (1 - e^{-\alpha \Delta t})
      + \sigma \int_t^{t+\Delta t} e^{-\alpha (t+\Delta t-u)} \, \mathrm{d}W_u \notag \\
    &= \tilde{C}_t e^{-\alpha \Delta t} + \mu (1 - e^{-\alpha \Delta t})
      + \sigma \int_t^{t+\Delta t} e^{-\alpha (t+\Delta t-u)} \, \mathrm{d}W_u
      \label{eq:ar1-recursion}
\end{align}
\end{subequations}
We have now shown that diffusion process is dependent on the previous emission $\tilde{C}_t$.\\ 
But for it to be an AR(1) process, we nedd to show that the drift term 
has mean 0, constant variance and is unoccorelated over time. [\textcolor{red}{MISSING REFERENCE}].\\
Let $\Delta t = c \in \mathbb{N}$ be a constant time step and given that $\alpha, t \in \mathbb{R}_{\geq 0}$, then from 
\ref{eq:ar1-recursion}, we can see that the drift term is given by
\begin{equation}
    \sigma \int_{t}^{t+\Delta t} e^{-\alpha (t+\Delta t-u)} \, \mathrm{d}W_u = 
\end{equation}
We recognize $e^{-\alpha (t+\Delta t-u)} $ as our stochastic term, which is determinitics, meaning that
that its ito intergratable as $\int_{0}^{T}e^{-2\alpha (t+\Delta t-u)} \, \mathrm{d}W_u < + \infty$. \\
Therefore, from \textcite[Thm.~1.9.1]{iacus2008},
We now know that the mean of the drit term is
\begin{equation}
    \mathbb{E}\left[\sigma \int_{t}^{t+\Delta t} e^{-\alpha (t+\Delta t-u)} \, \mathrm{d}W_u\right] = 0
\end{equation}
And that the variance is
\begin{equation}
    \text{Var}\left[\sigma \int_{t}^{t+\Delta t} e^{-\alpha (t+\Delta t-u)} \, \mathrm{d}W_u\right] = 
    \sigma^2 \int_{t}^{t+\Delta t} e^{-2\alpha (t+\Delta t-u)} \, \mathrm{d}u
\end{equation}
%Show that t to t + \delta T is 0 
%Show that the variance is independtent of t given Delta t
\begin{subequations}
\begin{align}
    \text{Var}\left[\sigma \int_{t}^{t+\Delta t} e^{-\alpha (t+\Delta t-u)} \, \mathrm{d}W_u\right]
    &= \sigma^2 \int_{t}^{t+\Delta t} e^{-2\alpha (t+\Delta t-u)} \, \mathrm{d}u \\
    &= \sigma^2 e^{-2\alpha (t+\Delta t)} \int_{t}^{t+\Delta t} e^{2\alpha u} \, \mathrm{d}u \\
    &= \sigma^2 e^{-2\alpha (t+\Delta t)} \left[\frac{1}{2\alpha} e^{2\alpha u}\right]_{t}^{t+\Delta t} \\
    &= \frac{\sigma^2}{2\alpha} e^{-2\alpha (t+\Delta t)} \left(e^{2\alpha (t+\Delta t)} - e^{2\alpha t}\right) \\
    &= \frac{\sigma^2}{2\alpha}\left(1 - e^{-2\alpha \Delta t}\right)
\end{align}
\end{subequations}
We see that the variance is independent of $t$ and 
that its constant over time. \\
We now only need to show that the covariance of the drift is 0. \\
We define that time steps as $t_1, t_2 \in \mathbb{R}_{\geq 0}$ such that $t_1 < t_2$ and that the time steps are disjoint, meaning that $[t_1, t_1 + \Delta t] \cap [t_2, t_2 + \Delta t] = \emptyset$ 
and that $t_1 + \Delta t = t_2$.
We then define $e_k$ as the drift term at time $t_k$ such that
\begin{equation}
\begin{aligned}
    \mathbb{E}\left[e_k\right] &= 0 \\ 
    \text{Var}(e_k) &= \sigma \int_{t_{k-1}}^{t_k} e^{-\alpha (t+\Delta t-u)} \, \mathrm{d}W_u = \frac{\sigma^2}{2\alpha}\left(1 - e^{-2\alpha \Delta t}\right)
\end{aligned}
\end{equation}

We now then see that the covariance of the drift term is given by, 
given that $j \neq k$
\begin{equation}
\begin{aligned}
      \text{Var}(e_k + e_j) &= \text{Var}(e_k) +  \text{Var}(e_j) + 2\text{Cov}(e_k, e_j) \\
      \Leftrightarrow 
\end{aligned}
\end{equation}



To compute $\text{Var}(e_k + e_j)$, we write both drift terms as integrals over
the same interval $[0, T]$, where $T \geq \max(t_k, t_j) + \Delta t$, by
introducing the deterministic functions
\begin{equation}
    f_k(u) = \mathbb{1}_{[t_k,\, t_k + \Delta t)}(u)\, e^{-\alpha (t_k + \Delta t - u)},
\end{equation}
such that
\begin{equation}
    e_k = \sigma \int_0^T f_k(u) \, \mathrm{d}W_u .
\end{equation}
Since $j \neq k$, the intervals $[t_k, t_k + \Delta t)$ and $[t_j, t_j + \Delta t)$
are disjoint, and therefore
\begin{equation}
    f_k(u) f_j(u) = 0 \quad \forall u \in [0, T].
\end{equation}
The sum $e_k + e_j = \sigma \int_0^T \left(f_k(u) + f_j(u)\right) \mathrm{d}W_u$
is again an Itô integral with a deterministic integrand, so by
\textcite[Thm.~1.9.1]{iacus2008}
\begin{equation}
\begin{aligned}
    \text{Var}(e_k + e_j)
    &= \sigma^2 \int_0^T \left(f_k(u) + f_j(u)\right)^2 \mathrm{d}u \\
    &= \sigma^2 \int_0^T f_k(u)^2 \, \mathrm{d}u
     + \sigma^2 \int_0^T f_j(u)^2 \, \mathrm{d}u
     + 2\sigma^2 \int_0^T f_k(u) f_j(u) \, \mathrm{d}u \\
    &= \text{Var}(e_k) + \text{Var}(e_j) + 0 .
\end{aligned}
\end{equation}
Comparing with
$\text{Var}(e_k + e_j) = \text{Var}(e_k) + \text{Var}(e_j) + 2\,\text{Cov}(e_k, e_j)$
gives
\begin{equation}
    \text{Cov}(e_k, e_j) = 0, \qquad j \neq k .
\end{equation}
Hence the drift terms have mean zero, constant variance and are uncorrelated,
and the discretized process is an AR(1) process.


\subsection{Modelling Production Rate}

By assuming that the prodction rate of the occupants is only depented on the 
amount of occupants in the room, we can model the production rate as a function of the
occupancy state of the room. Let $S(t)$ be the occupancy state of the room
at time $t$, then we can model the production rate as
\begin{equation}
    P(t) = \beta_{S(t)},
\end{equation}
where $\beta_{S(t)}$ is the production rate of the occupants in the room,
which is dependent on the occupancy state of the room. This gives a 
stepwise prodiuction rate function, which is only dependnt on the occupancy state of the room. 



%% --------- HMMM ----------
\section{Hidden Markov Model (HMM) theory} 
This section will present the underlying models and the theory behind it 
that models the diffusion process. Here to, inference algorithms and 
analysis methods will be presented. 

\subsection{The underlying theory}
%% Write some generel theory on HMM
Hidden markov models are markov chains where the underlying states are hidden and only a 
random variable that is dependent on the underlying state is observed. \\
This can be described as a markov chain with a set of hidden states $S = \{s_1, s_2, ..., s_N\}$ 
and a set of observations $O = \{o_1, o_2, ..., o_M\}$. 
The transition between the hidden states is governed by a transition matrix $\Gamma$ 
and the observations are generated by an emission matrix $B$. 
The initial state distribution is given by $\pi$. 
The HMM stucture can then debsrib the emissions as 
\begin{equation}
    P(O_t | S_t) = B_{S_t, O_t},
\end{equation}
where $B_{S_t, O_t}$ is the probability of observing $O_t$ given the hidden state $S_t$. 
The transition between the hidden states is given by
\begin{equation}
    P(S_t | S_{t-1}) = \Gamma_{S_{t-1}, S_t},
\end{equation}
where $\Gamma_{S_{t-1}, S_t}$ is the probability of transitioning from hidden state $S_{t-1}$ to hidden state $S_t$. 
The initial state distribution is given by
\begin{equation}
    P(S_1) = \pi_{S_1},
\end{equation}
where $\pi_{S_1}$ is the probability of starting in hidden state $S_1$.

\subsection{Likelihood}
Several methods can be used to estimate the parameters of a HMM, where at 
some of the more common methods are the Baum-Welch algorithm [\textcolor{red}{MISSING REFERENCE}].\\ 
In this project, we will however use the method of maximum likelihood estimation (MLE) 
to estimate the parameters of the HMM with the forward-filtering algorithm [\textcolor{red}{MISSING REFERENCE}]. \\

The forward algorithm implemented is called recursive filtering
\parencite{moller2024}, as it constructs the current state distribution
$p(X_t \mid Y_{1:t})$ recursively from the previous state distribution
$p(X_t \mid Y_{1:t-1})$.

\begin{algorithm}[H]
\caption{Recursive Filtering Forward Algorithm}
\label{alg:forward}
\begin{algorithmic}[1]
\Require Observations $\mathbf{y} = (y_1, \dots, y_T)$, covariates
    $\mathbf{x} = (x_1, \dots, x_T)$, initial state distribution
    $u_0 \in \mathbb{R}^{1 \times N}$, HMM parameters $\theta$
\Ensure Predicted distributions $\{u_t\}$, filtered distributions
    $\{u_{t|t}\}$, step likelihoods $\{f_t\}$
\For{$t = 1, \dots, T$}
    \State $\Gamma_t \gets \textsc{TransitionMatrix}(\theta, t, y_t, x_t)$
        \Comment{$N \times N$ transition matrix}
    \State $u_{t|t-1} \gets u_{t-1|t-1} \cdot \Gamma_t$
        \Comment{Predicted state distribution, $1 \times N$}
    \State $g_t \gets \textsc{Density}(\theta, t, y_t, x_t)$
        \Comment{Emission densities for all states, $1 \times N$}
    \State $f_t \gets \sum_{i=1}^{N} \left[ u_t \odot g_t \right]_i$
        \Comment{Marginal likelihood at time $t$}
    \State $f_t \gets \max(f_t, 10^{-10})$
        \Comment{Clip to avoid division by zero}
    \State $u_{t|t} \gets \dfrac{u_t \odot g_t}{f_t}$
        \Comment{Filtered state distribution (normalized)}
\EndFor
\State \Return $\{u_t\}$, $\{u_{t|t}\}$, $\{f_t\}$
\end{algorithmic}
\end{algorithm}

\noindent where $\odot$ denotes element-wise multiplication of vectors.
The log-likelihood is obtained by summing the log of the step likelihoods:
\begin{equation}
    \mathcal{L}(\theta) = \sum_{t=1}^{T} \log f_t .
\end{equation}
Each likelihood step $f_t$ is the sum of the predicted state distribution $u_t$ multiplied by the emission densities $g_t$ for all states, 
which gives the marginal likelihood of observing $y_t$ given the previous observations $y_{1:t-1}$.

\subsection{Forward Backward Filtering} 

\subsection{Backward Sampling} 
\subsection{Model Diagnostics}
